\documentclass[11pt]{article}
\usepackage[margin=1in]{geometry}
\usepackage{enumitem}
% =============================================================================
% Preambles/header.tex — shared LaTeX/Beamer preamble for this project
%
% Use from a Beamer lecture by prepending:
%     \input{header}
% and compiling with TEXINPUTS=../Preambles:$TEXINPUTS (the /compile-latex
% skill does this automatically).
%
% PALETTE CONTRACT. Color names defined here MUST match the names in
% Quarto/theme-template.scss so Beamer and Quarto renderings use the same
% palette. The scripts/check-palette-sync.sh script enforces this.
%
% To customize: edit the HEX values below AND the matching $variable in
% Quarto/theme-template.scss, then run scripts/check-palette-sync.sh.
% =============================================================================

% -----------------------------------------------------------------------------
% Engine detection + encoding
%
% Our /compile-latex skill uses XeLaTeX, where inputenc is unnecessary and
% can produce encoding warnings. Only load inputenc under pdfTeX.
% -----------------------------------------------------------------------------
\usepackage{iftex}
\ifPDFTeX
  \usepackage[utf8]{inputenc}
\fi

\usepackage{amsmath, amssymb, amsthm}
\usepackage{booktabs}
\usepackage{graphicx}

% -----------------------------------------------------------------------------
% PALETTE — mirrors Quarto/theme-template.scss variable names.
% Load xcolor and define colors BEFORE anything that references them
% (notably hyperref, hypersetup, and the Beamer color assignments).
% -----------------------------------------------------------------------------
\usepackage{xcolor}

\definecolor{primary-blue}{HTML}{012169}      % headings, accents
\definecolor{primary-gold}{HTML}{B9975B}      % emphasis, borders
\definecolor{highlight-yellow}{HTML}{F2A900}  % markers, alerts
\definecolor{light-bg}{HTML}{E8EDF5}          % title / section backgrounds
\definecolor{jet}{HTML}{1A1A1A}               % body text

% Semantic colors (used in diagrams and annotations)
\definecolor{positive}{HTML}{15803D}          % good / observed / identified
\definecolor{negative}{HTML}{B91C1C}          % bad / problematic / bias
\definecolor{neutral}{HTML}{525252}           % reference / context

% Highlight accents (match .hi-* classes in the SCSS)
\definecolor{hi-slate}{HTML}{314F4F}
\definecolor{hi-green}{HTML}{2E8B57}
\definecolor{hi-red}{HTML}{C41E3A}

% Semantic aliases so skill templates and snippets can write
% \textcolor{accent}{...} without hard-coding "primary-blue".
\colorlet{accent}{primary-blue}
\colorlet{accent2}{primary-gold}

% -----------------------------------------------------------------------------
% Hyperref — loaded AFTER xcolor + palette so link colors resolve.
% -----------------------------------------------------------------------------
\usepackage{hyperref}
\hypersetup{colorlinks=true, linkcolor=primary-blue, urlcolor=primary-blue,
            citecolor=primary-blue, pdfborder={0 0 0}}

% -----------------------------------------------------------------------------
% Course code — set once per deck (after \input{header}, before \begin{document})
% via \coursecode{CS401}. Shown in the footer on every slide so a deck's course
% is identifiable without opening the title page. Empty by default.
% -----------------------------------------------------------------------------
\makeatletter
\newcommand{\coursecode}[1]{\gdef\@coursecodetext{#1}}
\gdef\@coursecodetext{}
\makeatother

% -----------------------------------------------------------------------------
% Beamer theme assignments (only apply under Beamer).
%
% We detect Beamer via \@ifclassloaded{beamer}, which is the documented,
% non-internal way. This must be wrapped in \makeatletter/\makeatother
% because \@ifclassloaded contains an @-catcode character.
% -----------------------------------------------------------------------------
\makeatletter
\@ifclassloaded{beamer}{%
  \usefonttheme{professionalfonts}
  \setbeamercolor{structure}{fg=primary-blue}
  \setbeamercolor{frametitle}{fg=primary-blue}
  \setbeamercolor{title}{fg=primary-blue}
  \setbeamercolor{subtitle}{fg=primary-gold}
  \setbeamercolor{author}{fg=primary-blue}
  \setbeamercolor{institute}{fg=neutral}
  \setbeamercolor{date}{fg=primary-gold}
  \setbeamercolor{itemize item}{fg=primary-blue}
  \setbeamercolor{itemize subitem}{fg=primary-gold}
  \setbeamercolor{alerted text}{fg=primary-gold}
  \setbeamercolor{block title}{fg=white, bg=primary-blue}
  \setbeamercolor{block body}{bg=light-bg}
  % Semantic block variants: block=definition/concept (blue), exampleblock=
  % worked example (green/positive), alertblock=key takeaway or caution (gold).
  % Keep to <=2 colored boxes per slide across ALL three variants (INV-7).
  \setbeamercolor{block title example}{fg=white, bg=positive}
  \setbeamercolor{block body example}{bg=light-bg}
  \setbeamercolor{block title alerted}{fg=white, bg=primary-gold}
  \setbeamercolor{block body alerted}{bg=light-bg}

  % Minimal footer: course code (left) + page X / N (right), no navigation clutter
  \setbeamertemplate{navigation symbols}{}
  \setbeamertemplate{footline}{%
    \color{neutral}\footnotesize\hspace{0.7em}\@coursecodetext\hfill
    \insertframenumber/\inserttotalframenumber\hspace{0.7em}\vspace{0.3em}}
}{}
\makeatother

% -----------------------------------------------------------------------------
% TikZ setup — libraries the snippet gallery uses
% -----------------------------------------------------------------------------
\usepackage{tikz}
\usetikzlibrary{arrows.meta, positioning, calc, decorations.pathreplacing,
                fit, shapes.geometric, backgrounds}

% Shared TikZ styles — snippets and hand-written diagrams can pick these up
% via `\begin{tikzpicture}[every node/.style={...}]` or by naming specific
% styles (e.g., `\node[dag-node] (X) at (0,0) {X};`).
\tikzset{
  % Node styles for causal DAGs and similar
  dag-node/.style = {
    draw=primary-blue, fill=light-bg, rounded corners=2pt,
    minimum width=1.2cm, minimum height=0.9cm, align=center, font=\small
  },
  decision-node/.style = {
    draw=primary-gold, fill=white, diamond, aspect=1.8,
    minimum width=1.8cm, minimum height=1.2cm, align=center, font=\small,
    inner sep=1pt
  },
  % Edge styles — observed vs counterfactual
  observed-edge/.style = {-{Latex[length=2.5mm]}, thick, draw=jet},
  counterfactual-edge/.style = {-{Latex[length=2.5mm]}, thick, draw=neutral, dashed},
  confound-edge/.style = {-{Latex[length=2.5mm]}, thick, draw=negative, dashed},
  % Data-point styles for plots
  observed-dot/.style = {fill=primary-blue, draw=primary-blue, circle, inner sep=1.2pt},
  counterfactual-dot/.style = {fill=white, draw=neutral, circle, inner sep=1.2pt}
}

% -----------------------------------------------------------------------------
% Convenience macros
% -----------------------------------------------------------------------------
\newcommand{\muted}[1]{\textcolor{neutral}{#1}}
\newcommand{\key}[1]{\textcolor{primary-gold}{\textbf{#1}}}
\newcommand{\good}[1]{\textcolor{positive}{#1}}
\newcommand{\bad}[1]{\textcolor{negative}{#1}}

% Standout frame macro: full-bleed dark background for section transitions.
% Beamer-only (uses \begin{frame}/\setbeamercolor/\usebeamerfont) -- do not
% call from an article-class file (e.g. Notes/<CODE>/*-notes.tex). \muted,
% \key, \good, \bad, and \coursecode above are plain xcolor/gdef and are
% safe to use from either class; verified when Notes/article-class support
% was added.
% Expand: \transitionslide{Your transition title}
\newcommand{\transitionslide}[1]{%
  \begingroup
  \setbeamercolor{background canvas}{bg=primary-blue}
  \begin{frame}[plain]
    \centering
    \vfill
    {\usebeamerfont{title}\color{white}\Large #1\par}
    \vfill
  \end{frame}
  \endgroup
}

% Section divider frame (Beamer-only), an alternative to \transitionslide with a
% darker two-line style: near-black (jet) background, a small white label line
% above a larger gold title, and a thin gold rule along the bottom edge.
% Expand: \sectiondivider{Section 2}{Pointers}
\newcommand{\sectiondivider}[2]{%
  \begingroup
  \setbeamercolor{background canvas}{bg=jet}
  \begin{frame}[plain]
    \vspace*{3.0cm}
    {\color{white}\ttfamily\large #1\par}
    \vspace{0.5cm}
    {\color{primary-gold}\LARGE #2\par}
    \begin{tikzpicture}[remember picture, overlay]
      \draw[primary-gold, line width=2.5pt]
        ($(current page.south west)+(0,0.1)$) -- ($(current page.south east)+(0,0.1)$);
    \end{tikzpicture}
  \end{frame}
  \endgroup
}

% -----------------------------------------------------------------------------
% Channel watermark — "Concepts That Click" (YouTube).
%
% Article-class files (CompetitiveExam Q/A sets, Notes, Assignments) get a
% light diagonal draft watermark on every page plus a centered footer naming
% the channel. Beamer decks get a subtle TikZ background watermark instead
% (the footline already carries the course code). Centralized here so every
% MCQ PDF is branded without per-file edits.
% -----------------------------------------------------------------------------
\makeatletter
\@ifclassloaded{article}{%
  \usepackage{draftwatermark}
  \SetWatermarkText{Concepts That Click}
  \SetWatermarkScale{1}
  \SetWatermarkFontSize{2cm}
  \SetWatermarkLightness{0.86}
  \SetWatermarkAngle{45}
  \usepackage{fancyhdr}
  \pagestyle{fancy}
  \fancyhf{}
  \fancyfoot[C]{\footnotesize\textcolor{neutral}{Concepts That Click~|~YouTube: \href{https://www.youtube.com/@concepts.thatclick}{@concepts.thatclick}}}
  \renewcommand{\headrulewidth}{0pt}
  \renewcommand{\footrulewidth}{0pt}
  \fancypagestyle{plain}{%
    \fancyhf{}%
    \fancyfoot[C]{\footnotesize\textcolor{neutral}{Concepts That Click~|~YouTube: \href{https://www.youtube.com/@concepts.thatclick}{@concepts.thatclick}}}%
    \renewcommand{\headrulewidth}{0pt}%
    \renewcommand{\footrulewidth}{0pt}%
  }
}{}
\@ifclassloaded{beamer}{%
  \setbeamertemplate{background}{%
    \begin{tikzpicture}[remember picture, overlay]
      \node[rotate=45, opacity=0.055, align=center]
        at (current page.center)
        {\Huge\textcolor{neutral}{Concepts That Click}};
    \end{tikzpicture}%
  }
}{}
\makeatother 
\coursecode{JKSSB}
\renewcommand{\thesection}{28.\arabic{section}}
\newtheorem{example}{Example}[section]
\newenvironment{definitionbox}[1]{%
  \par\noindent\textbf{Definition (#1).}\ \itshape
}{\par\vspace{0.3em}}
\title{Access Database Basics --- Detailed Lecture Notes}
\author{JKSSB Computer Awareness}
\date{\today}

\begin{document}
\maketitle

\section{What a database is}
A \textbf{database} is an organised collection of related data, stored so
that it can be searched, updated, and retrieved easily. The data is not kept
as a loose pile of separate files; it is arranged in a planned structure.
A simple example is a school's information: the list of students, their
classes, and their marks all belong together as related information, and a
database keeps them that way.

A useful comparison is a spreadsheet. A spreadsheet (such as Excel) is built
for calculation on a simple grid. A database is built for storing and linking
large amounts of related data, and it can hold many tables that are related
to one another. When a record system grows large, a database keeps the data
consistent and avoids unnecessary duplication.

\begin{definitionbox}{Database}
An organised collection of related data, stored so that it can be searched,
updated, and retrieved easily.
\end{definitionbox}

\section{Database vs DBMS}
A database and a DBMS are not the same thing, and the difference is a common
exam point. The \textbf{database} is the collection of data itself. A
\textbf{DBMS} (Database Management System) is the software that manages that
data. The database is data; the DBMS is a program.

\begin{center}
\begin{tabular}{p{0.18\textwidth}p{0.35\textwidth}p{0.37\textwidth}}
\toprule
\textbf{Point} & \textbf{Database} & \textbf{DBMS} \\
\midrule
What it is & The organised collection of data & The software that manages the data \\
Nature & Data & Program \\
Example & A school's records & MS Access, Oracle, MySQL \\
\bottomrule
\end{tabular}
\end{center}

\begin{definitionbox}{DBMS (Database Management System)}
The software that stores a database and lets users insert, update, delete,
and retrieve data while controlling access and keeping the data consistent.
\end{definitionbox}

A \textbf{relational database management system} (RDBMS) is a DBMS that keeps
data in related tables. Microsoft Access is an example of an RDBMS.

\section{What a DBMS does}
A DBMS performs several jobs. It \emph{stores} the data and controls how it is
kept on disk. It lets users \emph{insert}, \emph{update}, \emph{delete}, and
\emph{retrieve} records. It controls \emph{access}, so that only permitted
users can read or change the data. It also keeps the data
\emph{consistent} --- different parts of the database are not allowed to
contradict each other --- and protects it from corruption.

Common DBMS software includes MS Access, Oracle, MySQL, and SQL Server. For
this course, the important point is that Access is a DBMS: it is the program
that manages a database, not the data itself.

\section{Microsoft Access in the Office suite}
\textbf{Microsoft Access} is the database program of the Microsoft Office
suite. In the suite, Word handles text documents, Excel handles spreadsheets,
PowerPoint handles presentations, and Access handles databases. Access is a
relational DBMS: it stores its data in tables, and it also holds the queries,
forms, and reports that work on that data. All of these objects are kept
inside a single database file for each database.

\begin{example}
A student is asked, ``Which MS Office program is used to store and manage
databases?'' The answer is \textbf{Access}. Excel is a spreadsheet program,
Word is a word processor, and PowerPoint is for presentations; none of them
is a DBMS.
\end{example}

\section{The objects in an Access database}
An Access database contains several kinds of object, each with one job.

\begin{center}
\begin{tabular}{p{0.16\textwidth}p{0.70\textwidth}}
\toprule
\textbf{Object} & \textbf{Job} \\
\midrule
Table & Stores the actual data in rows and columns \\
Query & Asks a question and retrieves selected data \\
Form & Provides a screen to enter and view data \\
Report & Presents data in a printable layout \\
\bottomrule
\end{tabular}
\end{center}

The data itself always lives in \textbf{tables}. The other objects work on
that data: a query retrieves it, a form is used to enter and view it, and a
report presents it for printing.

\section{Tables, records, and fields}
A \textbf{table} is the basic object where data is stored. Inside a table,
data is arranged in rows and columns. Each \textbf{column} is one
\textbf{field} --- one kind of information, such as Name or Marks. Each
\textbf{row} is one \textbf{record} --- one complete entry, such as the full
details of one student. A table therefore holds many records, and every
record has the same set of fields.

\begin{center}
\begin{tabular}{p{0.16\textwidth}p{0.34\textwidth}p{0.30\textwidth}}
\toprule
\textbf{Term} & \textbf{Meaning} & \textbf{Position in a table} \\
\midrule
Field & One kind of information & A column \\
Record & One complete entry & A row \\
Table & A collection of related records & The whole grid \\
\bottomrule
\end{tabular}
\end{center}

\begin{definitionbox}{Record}
One complete entry in a table, held in a single row and made up of one value
for each field.
\end{definitionbox}

\begin{definitionbox}{Field}
One kind of information in a table, held in a single column; each field has a
name and a data type.
\end{definitionbox}

\begin{example}
In a Students table the columns are RollNo, Name, Marks, and Passed. One row
--- for example, ``1, Ayaan, 78, Yes'' --- is one record. The Name column is
the Name field. The whole set of rows together is the table.
\end{example}

\section{Data types in Access}
Every field has a \textbf{data type}, which is the kind of value it can hold.
Common Access data types are:

\begin{center}
\begin{tabular}{p{0.22\textwidth}p{0.60\textwidth}}
\toprule
\textbf{Data type} & \textbf{Used for} \\
\midrule
Short Text & Names, codes, and other text \\
Number & Values used in calculations \\
Date/Time & Dates and times \\
Yes/No & Only True or False (a Boolean value) \\
AutoNumber & A unique number added automatically for each new record \\
\bottomrule
\end{tabular}
\end{center}

Because a field accepts only values that match its data type, the data stays
clean. A field set to Number will not accept ordinary text, and a Yes/No
field accepts only a true-or-false value.

\section{Field properties and validation}
A \textbf{field property} is a setting that describes how a field behaves.
Some important properties are \textbf{Field Size}, which limits how much a
field can hold; \textbf{Default Value}, which fills in a value
automatically; \textbf{Required}, which, when set to Yes, stops the field
from being left blank; and \textbf{Input Mask}, which fixes the pattern of
entry, such as the format of a phone number.

\textbf{Validation} means checking that entered data is sensible and allowed.
A \textbf{Validation Rule} states the condition a value must meet. For
example, the rule \texttt{Between 0 And 100} accepts only values in that
range. A \textbf{Validation Text} is the message shown when the rule is
violated.

\begin{definitionbox}{Validation}
The checking of entered data against a rule so that only sensible, allowed
values are stored in a field.
\end{definitionbox}

\begin{example}
On a Marks field, the Validation Rule \texttt{Between 0 And 100} is set. A
user types 150. Access rejects the entry because it fails the rule, and shows
the Validation Text as a message. A valid value such as 78 is accepted
normally.
\end{example}

\section{Primary key and index}
A \textbf{primary key} is a field that uniquely identifies each record. Its
value is different for every record and is never left blank. In the Students
table, RollNo can serve as the primary key because no two students share a
roll number. An AutoNumber field is often used as a ready-made primary key,
because Access fills it with a fresh unique number for every new record.

An \textbf{index} is a special structure that makes searching and sorting a
table faster. It works like the index of a book: instead of reading every
page, it points straight to the right place. A primary key is automatically
indexed. An index takes extra space and can slightly slow down the work of
adding data, so it is created only where it is genuinely useful.

\begin{definitionbox}{Primary key}
A field, or set of fields, whose value uniquely identifies every record in a
table; it cannot contain duplicates and cannot be left blank.
\end{definitionbox}

\section{Exam retrieval and common confusions}
Five distinctions prevent most errors on this topic.
\begin{enumerate}
  \item A \textbf{database} is the data; a \textbf{DBMS} is the software that
    manages it. Do not swap them.
  \item A \textbf{field} is a column; a \textbf{record} is a row; a
    \textbf{table} is the whole collection of related records.
  \item \textbf{Validation} checks the value entered; an \textbf{index}
    speeds up searching. They do different jobs.
  \item A \textbf{primary key} is unique and cannot be blank; it cannot
    contain duplicate values.
  \item Microsoft \textbf{Access} is a DBMS and part of the MS Office suite;
    it is not a spreadsheet.
\end{enumerate}

\begin{example}
An item asks which object stores the actual data in an Access database.
Trace the objects: a query retrieves data, a form is for entry, and a report
is for printing, but only the \textbf{table} stores the data. If the stem
instead asks which field uniquely identifies a record, the answer is the
\textbf{primary key}; if it asks which field setting stops a blank entry, the
answer is \textbf{Required}.
\end{example}

\subsection*{Exercises}
\begin{enumerate}[label=\arabic*.]
  \item Define a database and explain how it differs from a DBMS.
  \item Distinguish a field, a record, and a table, stating the position of
    each in a table.
  \item Name any four Access data types and give one use of each.
  \item Explain the difference between a validation rule and an index.
  \item What is a primary key, and why can it not contain duplicate values?
\end{enumerate}

\subsection*{Solutions}
\begingroup\small
\begin{enumerate}[label=\arabic*.]
  \item A database is an organised collection of related data. A DBMS is the
    software that manages that data --- it stores it, allows insert, update,
    delete, and retrieve operations, and controls access. The database is the
    data; the DBMS is the program.
  \item A field is one kind of information and is a column in a table. A
    record is one complete entry and is a row in a table. A table is the
    whole collection of related records, arranged in rows and columns.
  \item Any four of: Short Text (names and codes), Number (values used in
    calculations), Date/Time (dates and times), Yes/No (True or False),
    AutoNumber (a unique number for each new record).
  \item A validation rule checks that entered data meets a condition, such as
    \texttt{Between 0 And 100}, and rejects values that fail it. An index is a
    structure that speeds up searching and sorting; it does not check values.
  \item A primary key is a field whose value uniquely identifies each record.
    It cannot contain duplicates because two records with the same key value
    could not be told apart, and uniqueness is the whole purpose of the key.
\end{enumerate}
\endgroup
\end{document}
